\section{The contribution of short flags on the moderate region}

% BEGIN C83N5
\paragraph{C83N5. Counting nearby isometries on each arithmetic flag.}
Keep the point lattice and metrics of C83N3. Write $S(t)$ for the step
of its canonical filtration containing the slopes of normalized size
at most $t$. We use the canonical filtration theorem of Grayson,
\emph{Reduction theory using semistability}, Theorem~1.18 and
Definition~1.19. It applies to a projective integer lattice with the
Hermitian norms at its complex places. Changing log covolume by a
linear function of the rank and taking the positive constant multiple
in C83N3 leaves its vertices fixed. Thus its graded lattices are
semistable with strictly increasing sizes. Every nonzero sublattice
of a grade has size at least its size, and every nonzero torsion-free
quotient has size at most its size, by multiplicativity of covolume.
The same polygon gives the corresponding least and largest bounds on a
consecutive group of grades. Dualizing reverses the order and inverts
the sizes by C83N3 and uniqueness of the semistable filtration.

Let $\tau_0<\cdots<\tau_{a-1}<1/2$ be all the short sizes on a point
of the original type $j=(j_i)$, of ranks $j_i$, and $J=\sum_i j_i$.
Group the sizes in $[1/2,2]$ in one middle grade. The reflected
grades have sizes $1/\tau_i$ in reverse order. Denote the resulting
steps by $G_0=0,\ldots,G_{2a+1}=\Lambda_x$ and the grades by
$V_p=G_{p+1}/G_p$. Self-duality gives
$G_i^\perp=G_{2a+1-i}$ inside $\Lambda_x$. In particular the short
subspace $G_a$ is isotropic, so $J\le\min(p_+,q)$ at a real place
of signature $(p_+,q)$. The middle rank is $n-2J$, possibly zero.

For $0<\delta<1/4$, take any rational unitary $g$ in the original
near set, so $g$ preserves $\Lambda_x$ and $\|g-1\|_{\rm HS}\le\delta$
at each complex place. Then $g-1$ sends $S(t)$ into $S(\delta t)$.
Indeed otherwise the map from $S(t)$ to the quotient by
$S(\delta t)$ is nonzero. Its quotient source has size at most
$t$, and its image has size strictly greater than $\delta t$.
All induced operator norms are at most $\delta$, contradicting the
determinant inequality of C83N3. It follows that $g-1$ strictly lowers
the grouped filtration, also on the middle grade since $2\delta<1/2$.

For $k\ge1$, two such elements agree through the preceding grades of
endomorphisms when $A=g-g'$ lowers the filtration at least $k$ steps.
Here $G_i=0$ for $i<0$. In one such equivalence class, subtract a
fixed representative. Its degree-$k$ blocks are integral maps
$V_p\to V_l$, $l=p-k$. Equality of all these blocks is exactly
agreement up to a drop of $k+1$. Pair $V_l$ with $V_{2a-l}$ by the
Hermitian form. Unitarity gives
\[
            h(Ax,gy)+h(g'x,Ay)=0 .
\]
For $x\in G_{p+1}$, $y\in G_{p'+1}$ and $p+p'=2a+k$, the terms
involving $(g-1)y$ and $(g'-1)x$ vanish by orthogonality of the
steps, since they lower one extra step. Thus the degree-$k$ blocks
are negatives of their reflected adjoints. It suffices to record
the pairs $(l,p)$ with $l<p$, $l+p\le2a$ and $l<a$. Those with
$l+p<2a$ have real dimension $2\dim V_l\dim V_p$ per complex place;
on $l+p=2a$ this is a skew-Hermitian form on $V_p$, of real
dimension $(\dim V_p)^2$. These dimensions follow by choosing
bases on paired grades. The norms on the block spaces are those
induced by the fixed positive metrics on the graded pieces.

Every recorded block differs from the chosen base block by norm at
most $2\delta$ at each place, since it arises by restriction and
orthogonal quotient projection. Distinct values in the same block
are separated in geometric mean norm by at least $t_l/t_p$, where
$t_l=\tau_l$ and $t_p$ is the corresponding size or $2$ on the
middle grade. In fact the difference is integral on the two
quotient lattices. Its nonzero image has size at least $t_l$, its
source quotient at most $t_p$, and C83N3's determinant argument
applies. Its maximum norm over the places is therefore at least
$t_l/t_p$. Packing that block as in C83N3 gives at most
$(1+4\delta t_p/t_l)^{fD}$ possibilities if its real dimension at
one place is $D$. Multiplying these bounds for all independent
blocks of each drop and then the finite successive drops counts
the whole class. After drop $2a$ agreement on every block implies
equality because $\Lambda_x$ spans its rational space.
The empty short flag gives the identity alone by C83N3.

Put $H(z)=1+12\delta/z$. The short--short, short--reflected,
short--middle and self-paired blocks respectively give the bound
\[
 \begin{aligned}
 P_\delta(j,\tau)={}&
 \prod_{i<l}
    \bigl[H(\tau_i/\tau_l)H(\tau_i\tau_l)\bigr]^{2j_i j_l f}\\
 &\times\prod_i H(\tau_i)^{2j_i(n-2J)f}
                             H(\tau_i^2)^{j_i^2f}.
 \end{aligned}
\]
For the middle the raw numerator is $8\delta$ and all other raw
numerators are $4\delta$, both bounded by $12\delta$. This proves
finiteness and the native cardinality bound on the flag of the actual
point. The count includes the identity and applies before any further
restriction to the individual cylinder support.
% END C83N5

% BEGIN C83N6
\paragraph{C83N6. The two flag estimates sufficient for moderate near terms.}
Use the radius and the counting hypotheses of C83N2 and C83N4, with
$1\le q$, and let $L=nf\log n$, $X=nf\,\mathrm{rd}\,g b$.
Retain the two native inputs $\mathrm{FlagMeasureBound}$ and
$\mathrm{ConstantBound}$. In their application to a type of short
grades of total rank $1\le J\le q$, the pushforward of the measure
on that part of the fundamental domain under its size tuple is
dominated by
\[
 \mu(\mathscr F)\,\kappa_{p,q,j}\,
 P_{\rm id}(j,\tau)\prod_i(2fj_i)\,\frac{d\tau_i}{\tau_i},
 \qquad \kappa_{p,q,j}\le e^{-JL/2},
\]
for sufficiently large $n$. In fact $q^2\le n$ implies $2q+2\le n$ eventually, and $\log n\ge15$ eventually. The row identities $p+q=n$, $2v_{\rm inv}(E)=\mathrm{rd}$ and $f_{\rm deg}(E)=f$ give the discriminant and degree guards in the second input from the counting hypotheses. Here
$P_{\rm id}$ is C83N5's product with $H(z)=z$.
The measure assertion on measurable subsets suffices for
nonnegative functions of $\tau$ by simple-function approximation.

At a point in the support the kernel is bounded by
$a_0=(\prod_v d(\sigma_v))\operatorname{Re}e(1)$.
The normalized archimedean coefficients are coefficients of a unit
vector in a unitary representation. The finite self-adjoint
convolution idempotent satisfies $e=e*e^*$, so its value at one is its
squared $L^2$ norm and the Cauchy--Schwarz inequality applied to
each translate gives $|e(g)|\le e(1)$. By C83N2 each contributing
near element satisfies the norm bound in C83N5. The nonidentity
near sum vanishes at least size $\ge1/2$ by C83N3. At each other
moderate point it is at most $a_0P_\delta(j,\tau)$ on its unique
short flag, all of whose sizes lie between $X^{-11q}$ and $1/2$.

Every argument $z$ of $H$ in C83N5 is at most one. Multiplying each
factor by its density factor gives $z+12\delta\le e^{12\delta}$.
The sum of the exponents is at most $3Jnf$, by summation of the
displayed exponents and $2J\le n$. Thus
$P_\delta P_{\rm id}\le e^{36\delta qnf}\le2$ by C83N2.
Integrate over the enclosing product of intervals; each factor
of the logarithmic measure has mass at most
$Y=22f q^2\log X$. With $D^\star=a_0\mu(\mathscr F)$ the
contribution of each composition is at most $2D^\star e^{-JL/2}Y^a$.

For fixed $C_g$, uniformly in the counting range, $Y\le e^{L/8}$
for large $n$. Indeed $q^2\le n$, $\log\mathrm{rd}\le e^n$,
$\log b\le e^{L/16}$, $\log f\le\sqrt n$, and $g\le b^{C_g}$.
Eventually $n\le L/16$, so
$\log X\le(C_g+4)e^{L/16}$ and
$22nf(C_g+4)\le e^{L/16}$.
There are $2^{J-1}$ ordered compositions of $J$ as positive
parts and $a\le J$ for each. Their total number for $1\le J\le q$
is less than $2^q$. Since $2^{q+1}\le e^{L/8}$ eventually, the
integral on the entire moderate region is at most
$D^\star e^{-L/4}$. This is the original $\mathrm{NearModerate}$
for each finite row projector. It requires the stated quantitative
flag measure and constant. The local coefficient decay in C83N4
enters the far-term integral separately.
% END C83N6
